\clearpage
\fancyhead[L]{\small HARDY’S PENDULUM CLOCK.}
\hypertarget{native-458-hardy}{}
\begin{center}{\large HARDY’S PENDULUM CLOCK.}\end{center}

\begin{center}
\rule{0.55\linewidth}{0.4pt}\\[2pt]
Göttingische gelehrte Anzeigen. No. 101. pp. 1001\,..\,1002. 26 June 1826.\\[2pt]
\rule{0.55\linewidth}{0.4pt}
\end{center}

His Royal Highness the Duke of \textsc{Sussex} has again shown his gracious goodwill to the university by giving the observatory a valuable gift. It is an astronomical pendulum clock by \textsc{Hardy}. For some months it has already served as the main clock for meridian observations. In that service it has fully lived up to the expectations raised by the maker’s name and the excellence of his work. We give here an account of its rate so far, whose uniformity leaves nothing to be desired. The compensation is achieved in \SourceLeaf{459}the simplest possible way. At its lower end, the pendulum rod carries, in place of the usual lenticular bob, a frame holding a cylindrical glass vessel $6\frac{1}{2}$ inches high and 2 inches across, nearly full of mercury. If the compensation is found to be too weak, some mercury must be added. In the opposite case, some must be taken out. In this way the compensation can be adjusted very finely once the clock’s rate has been observed at the extremes of temperature. The following table shows that as the temperature rises, the daily rate decreases quite gradually by half a second, and that the compensation is therefore still somewhat too weak.

\begin{center}
\emph{Offset and rate of \textup{\textsc{Hardy}}’s regulator relative to sidereal time, reduced to noon.}
\end{center}

\begin{center}
\begingroup\small\setlength{\tabcolsep}{5pt}
\begin{tabular}{lr|r|r}
 & & Offset & Daily rate \\
\hline
1826 & March 9 & $+0^{\mathrm s}83$ & $+0.42$ \\
 & 11 & $1.67$ & $+0.52$ \\
 & 14 & $3.22$ & $+0.45$ \\
 & 19 & $5.46$ & $+0.35$ \\
 & April 20 & $16.55$ & $+0.14$ \\
 & 21 & $16.69$ & $+0.07$ \\
 & 22 & $16.76$ & $+0.21$ \\
 & May 2 & $19.09$ & $+0.16$ \\
 & 8 & $20.05$ & $+0.18$ \\
 & 9 & $20.23$ & $+0.25$ \\
 & 11 & $20.74$ & $+0.01$ \\
\end{tabular}
\par\medskip
\begin{tabular}{lr|r|r}
 & & Offset & Daily rate \\
\hline
1826 & May 14 & $+20^{\mathrm s}78$ & $+0.08$ \\
 & 16 & $20.94$ & $+0.21$ \\
 & 18 & $21.37$ & $+0.15$ \\
 & 20 & $21.68$ & $+0.02$ \\
 & 23 & $21.73$ & $+0.10$ \\
 & 26 & $22.03$ & $-0.06$ \\
 & 30 & $21.78$ & $-0.18$ \\
 & June 11 & $19.57$ & $-0.22$ \\
 & 16 & $18.45$ & $-0.18$ \\
 & 17 & $18.27$ &  \\
\end{tabular}
\endgroup
\end{center}

